\SourceLeaf{153}probability theory. I can hardly agree, though, to call what one obtains in this way the \textls{most probable} value. In fact, finding the most probable value is a mathematically quite definite problem, but by its nature it presupposes knowledge of the law of errors. It leads to arithmetic means \textls{only} in the single case in which that law is expressed in the form $e^{-kkxx}$. Generally speaking, the probability of the most probable value is itself only infinitely small. If $a$ is the most probable value in a \textls{continuous} set, this really means only that

\Needspace{7\baselineskip}
\begin{quote}
the probability that the true value lies between $a-\omega$ and $a+\omega$ is greater than the probability that it lies between any other pair of limits just as far apart, provided that $\omega$ is infinitely small.
\end{quote}

On closer inspection, however, precisely for this reason such a most probable value has little practical interest, much less than the value for which the error to be feared is least harmful on average. This is why, apart from other reasons that are of course just as important or much more so, I have preferred the second principle, which must not be confused with the first.

. . . . . . . . . . . .

\Needspace{6\baselineskip}
\begin{center}\textsc{Gauss} to \textsc{Bessel}. Göttingen, 28 February 1839.\end{center}

. . . . . . . . . . . . . . . .

I read with great interest your paper in the Astronomische Nachrichten on how the law governing the probability of observational errors arising from combined sources approaches the formula $e^{-\frac{xx}{hh}}$. To speak frankly, my interest was less in the matter itself than in your presentation. I have been familiar with the matter for many years, though I myself never got around to working out the development in full.

